% Alur Algoritma Euclidean dan pemakaiannya untuk mencari invers modulo.
% Dirujuk dari section-01.tex dengan % Alur Algoritma Euclidean dan pemakaiannya untuk mencari invers modulo.
% Dirujuk dari section-01.tex dengan % Alur Algoritma Euclidean dan pemakaiannya untuk mencari invers modulo.
% Dirujuk dari section-01.tex dengan % Alur Algoritma Euclidean dan pemakaiannya untuk mencari invers modulo.
% Dirujuk dari section-01.tex dengan \input{figures/alur-euclid-diperluas}.
\begin{figure}[htbp]
    \centering
    \begin{tikzpicture}[
        kotak/.style={draw, rounded corners, align=center, font=\small,
                      inner sep=5pt, minimum width=7.2cm},
        panah/.style={-{Latex[length=2.4mm]}, thick},
    ]
        \node[kotak, fill=gray!10] (masuk) at (0,0)
            {Masukan $a$ dan $n$ dengan $0 < a < n$};
        \node[kotak, fill=orange!12] (bagi) at (0,-1.5)
            {$r \leftarrow a \bmod n$; \quad $a \leftarrow n$; \quad $n \leftarrow r$};
        \node[kotak, fill=orange!12] (cek) at (0,-3.0)
            {Apakah $r = 0$?};
        \node[kotak, fill=green!10] (pbb) at (0,-4.5)
            {$\text{PBB}(a_0, n_0) = a$};
        \node[kotak, fill=blue!6] (invers) at (0,-6.0)
            {Substitusi balik: peroleh $x$ dengan $a_0 x \equiv 1 \pmod{n_0}$};

        \draw[panah] (masuk) -- (bagi);
        \draw[panah] (bagi) -- (cek);
        \draw[panah] (cek) -- node[right, font=\scriptsize, fill=white, inner sep=1pt]
            {ya} (pbb);
        \draw[panah] (pbb) -- node[right, font=\scriptsize, fill=white, inner sep=1pt]
            {PBB $= 1$} (invers);

        % Simpai ulang selama sisa bagi belum nol
        \draw[panah] (cek.east) -- (4.4,-3.0) -- (4.4,-1.5) -- (bagi.east)
            node[midway, above, font=\scriptsize, fill=white, inner sep=1pt] {belum};

        % Cabang kegagalan: PBB bukan 1
        \node[kotak, fill=red!6] (gagal) at (0,-7.5)
            {Jika $\text{PBB}(a_0, n_0) \neq 1$: invers modulo tidak ada};
        \draw[panah] (invers) -- (gagal);
    \end{tikzpicture}
    \caption{Alur pencarian invers modulo: Algoritma Euclidean mengulang pembagian sampai sisa bagi nol untuk memperoleh PBB, lalu substitusi balik menghasilkan $x$ sedemikian sehingga $a x \equiv 1 \pmod{n}$}
    \label{fig:alur-euclid-diperluas}
\end{figure}
.
\begin{figure}[htbp]
    \centering
    \begin{tikzpicture}[
        kotak/.style={draw, rounded corners, align=center, font=\small,
                      inner sep=5pt, minimum width=7.2cm},
        panah/.style={-{Latex[length=2.4mm]}, thick},
    ]
        \node[kotak, fill=gray!10] (masuk) at (0,0)
            {Masukan $a$ dan $n$ dengan $0 < a < n$};
        \node[kotak, fill=orange!12] (bagi) at (0,-1.5)
            {$r \leftarrow a \bmod n$; \quad $a \leftarrow n$; \quad $n \leftarrow r$};
        \node[kotak, fill=orange!12] (cek) at (0,-3.0)
            {Apakah $r = 0$?};
        \node[kotak, fill=green!10] (pbb) at (0,-4.5)
            {$\text{PBB}(a_0, n_0) = a$};
        \node[kotak, fill=blue!6] (invers) at (0,-6.0)
            {Substitusi balik: peroleh $x$ dengan $a_0 x \equiv 1 \pmod{n_0}$};

        \draw[panah] (masuk) -- (bagi);
        \draw[panah] (bagi) -- (cek);
        \draw[panah] (cek) -- node[right, font=\scriptsize, fill=white, inner sep=1pt]
            {ya} (pbb);
        \draw[panah] (pbb) -- node[right, font=\scriptsize, fill=white, inner sep=1pt]
            {PBB $= 1$} (invers);

        % Simpai ulang selama sisa bagi belum nol
        \draw[panah] (cek.east) -- (4.4,-3.0) -- (4.4,-1.5) -- (bagi.east)
            node[midway, above, font=\scriptsize, fill=white, inner sep=1pt] {belum};

        % Cabang kegagalan: PBB bukan 1
        \node[kotak, fill=red!6] (gagal) at (0,-7.5)
            {Jika $\text{PBB}(a_0, n_0) \neq 1$: invers modulo tidak ada};
        \draw[panah] (invers) -- (gagal);
    \end{tikzpicture}
    \caption{Alur pencarian invers modulo: Algoritma Euclidean mengulang pembagian sampai sisa bagi nol untuk memperoleh PBB, lalu substitusi balik menghasilkan $x$ sedemikian sehingga $a x \equiv 1 \pmod{n}$}
    \label{fig:alur-euclid-diperluas}
\end{figure}
.
\begin{figure}[htbp]
    \centering
    \begin{tikzpicture}[
        kotak/.style={draw, rounded corners, align=center, font=\small,
                      inner sep=5pt, minimum width=7.2cm},
        panah/.style={-{Latex[length=2.4mm]}, thick},
    ]
        \node[kotak, fill=gray!10] (masuk) at (0,0)
            {Masukan $a$ dan $n$ dengan $0 < a < n$};
        \node[kotak, fill=orange!12] (bagi) at (0,-1.5)
            {$r \leftarrow a \bmod n$; \quad $a \leftarrow n$; \quad $n \leftarrow r$};
        \node[kotak, fill=orange!12] (cek) at (0,-3.0)
            {Apakah $r = 0$?};
        \node[kotak, fill=green!10] (pbb) at (0,-4.5)
            {$\text{PBB}(a_0, n_0) = a$};
        \node[kotak, fill=blue!6] (invers) at (0,-6.0)
            {Substitusi balik: peroleh $x$ dengan $a_0 x \equiv 1 \pmod{n_0}$};

        \draw[panah] (masuk) -- (bagi);
        \draw[panah] (bagi) -- (cek);
        \draw[panah] (cek) -- node[right, font=\scriptsize, fill=white, inner sep=1pt]
            {ya} (pbb);
        \draw[panah] (pbb) -- node[right, font=\scriptsize, fill=white, inner sep=1pt]
            {PBB $= 1$} (invers);

        % Simpai ulang selama sisa bagi belum nol
        \draw[panah] (cek.east) -- (4.4,-3.0) -- (4.4,-1.5) -- (bagi.east)
            node[midway, above, font=\scriptsize, fill=white, inner sep=1pt] {belum};

        % Cabang kegagalan: PBB bukan 1
        \node[kotak, fill=red!6] (gagal) at (0,-7.5)
            {Jika $\text{PBB}(a_0, n_0) \neq 1$: invers modulo tidak ada};
        \draw[panah] (invers) -- (gagal);
    \end{tikzpicture}
    \caption{Alur pencarian invers modulo: Algoritma Euclidean mengulang pembagian sampai sisa bagi nol untuk memperoleh PBB, lalu substitusi balik menghasilkan $x$ sedemikian sehingga $a x \equiv 1 \pmod{n}$}
    \label{fig:alur-euclid-diperluas}
\end{figure}
.
\begin{figure}[htbp]
    \centering
    \begin{tikzpicture}[
        kotak/.style={draw, rounded corners, align=center, font=\small,
                      inner sep=5pt, minimum width=7.2cm},
        panah/.style={-{Latex[length=2.4mm]}, thick},
    ]
        \node[kotak, fill=gray!10] (masuk) at (0,0)
            {Masukan $a$ dan $n$ dengan $0 < a < n$};
        \node[kotak, fill=orange!12] (bagi) at (0,-1.5)
            {$r \leftarrow a \bmod n$; \quad $a \leftarrow n$; \quad $n \leftarrow r$};
        \node[kotak, fill=orange!12] (cek) at (0,-3.0)
            {Apakah $r = 0$?};
        \node[kotak, fill=green!10] (pbb) at (0,-4.5)
            {$\text{PBB}(a_0, n_0) = a$};
        \node[kotak, fill=blue!6] (invers) at (0,-6.0)
            {Substitusi balik: peroleh $x$ dengan $a_0 x \equiv 1 \pmod{n_0}$};

        \draw[panah] (masuk) -- (bagi);
        \draw[panah] (bagi) -- (cek);
        \draw[panah] (cek) -- node[right, font=\scriptsize, fill=white, inner sep=1pt]
            {ya} (pbb);
        \draw[panah] (pbb) -- node[right, font=\scriptsize, fill=white, inner sep=1pt]
            {PBB $= 1$} (invers);

        % Simpai ulang selama sisa bagi belum nol
        \draw[panah] (cek.east) -- (4.4,-3.0) -- (4.4,-1.5) -- (bagi.east)
            node[midway, above, font=\scriptsize, fill=white, inner sep=1pt] {belum};

        % Cabang kegagalan: PBB bukan 1
        \node[kotak, fill=red!6] (gagal) at (0,-7.5)
            {Jika $\text{PBB}(a_0, n_0) \neq 1$: invers modulo tidak ada};
        \draw[panah] (invers) -- (gagal);
    \end{tikzpicture}
    \caption{Alur pencarian invers modulo: Algoritma Euclidean mengulang pembagian sampai sisa bagi nol untuk memperoleh PBB, lalu substitusi balik menghasilkan $x$ sedemikian sehingga $a x \equiv 1 \pmod{n}$}
    \label{fig:alur-euclid-diperluas}
\end{figure}
